% Typografia wzorowana na „Miara i całka” G. Plebanka:
% PL Roman, 12 pt, A4, klasyczne rozdziały i podrozdziały.
\usepackage[utf8]{inputenc}
\usepackage[OT4,plmath]{polski}
\usepackage{amsmath,amssymb,amsthm}
\usepackage[left=30.56mm,right=26.2mm,top=30mm,bottom=28mm]{geometry}
\usepackage{microtype}
\usepackage{graphicx}
\usepackage{enumitem}
\usepackage[hidelinks,unicode]{hyperref}

\hypersetup{
  pdftitle={Teoria miary},
  pdfsubject={Skrypt z wykładów},
  linktoc=all,
  bookmarksnumbered=true
}
\graphicspath{{../grafika/}}
\setcounter{tocdepth}{1}
\setcounter{secnumdepth}{2}
\setlength{\parindent}{1.5em}
\setlength{\parskip}{0pt}
\setlist[enumerate]{label=\arabic*),leftmargin=2.1em,itemsep=0.25em}
\allowdisplaybreaks[1]

\theoremstyle{definition}
\newtheorem{definicja}{Definicja}[chapter]
\newtheorem{przyklad}[definicja]{Przykład}
\newtheorem{uwaga}[definicja]{Uwaga}
\theoremstyle{plain}
\newtheorem{fakt}[definicja]{Fakt}
\renewcommand{\proofname}{Dowód}

\newcommand{\R}{\mathbb{R}}
\newcommand{\N}{\mathbb{N}}
\newcommand{\A}{\mathcal{A}}
\newcommand{\Rbar}{\overline{\mathbb{R}}}

